% ============ 21 比较定理 (原书 pp.130-134) ============
\chapter{比较定理}\label{sec:21}
\index{比较定理}
\index{Riemann 存在定理}
\index{初等纤维化}
\index{解析簇}
\index{GAGA}

设 $X$ 是非奇异簇.则 $X$ 可以以自然的方式赋予复流形的结构.我把 $X$ 看作复流形时写为 $X^{\mathrm{an}}$，把 $X$ 看作带复拓扑的拓扑空间时写为 $X^{\mathrm{cx}}$($X^{\mathrm{an}}$ 就是 $X^{\mathrm{cx}}$ 连同环层).

\begin{theorem}\label{thm:21.1}
设 $X$ 是 $\mathbb{C}$ 上的非奇异簇.对任何有限 Abel 群 $\Lambda$ 与 $r \geq 0$，$H^{r}(X_{\mathrm{et}}, \Lambda) \cong H^{r}(X^{\mathrm{cx}}, \Lambda)$.
\end{theorem}

\begin{remark}\label{rem:21.2}
定理对奇异簇也成立，但这时陈述稍困难(需要知道解析簇)，证明也稍长.定理对一切可构造层也成立.见 SGA 4, XVI 4.

在两种拓扑中，$H^{0}(X, \Lambda) = \Lambda^{\pi_0(X)}$，其中 $\pi_0(X)$ 是 $X$ 的连通分支集.于是，对 $r = 0$，定理无非断言 $X$ 对 Zariski 拓扑的连通分支集与 $X^{\mathrm{cx}}$ 对复拓扑的连通分支集相同，或者等价地：若 $X$ 对 Zariski 拓扑连通，则它对复拓扑连通.

这是有点令人惊讶的结果.设 $X = \mathbb{A}^{1}$.当然，它对两种拓扑都连通(Zariski 拓扑的非空开子集是只省略有限多个点的那些；复拓扑使 $X$ 同胚于 $\mathbb{R}^{2}$).当我们从 $X$ 中去掉一个圆时，它对复拓扑不连通，但对 Zariski 拓扑仍连通.这并不与定理矛盾，因为去掉圆的 $\mathbb{A}^{1}_{\mathbb{C}}$ 不是代数簇.

设 $X$ 是连通非奇异(从而不可约)曲线.我们证明它对复拓扑连通.给 $X$ 去掉或加上有限个点不会改变它对复拓扑是否连通，因此可以假定 $X$ 射影.设 $X$ 是两个非空开(从而闭)集 $X_1$ 与 $X_2$ 的不交并.由 Riemann-Roch 定理，存在只在 $X_1$ 中有极点的 $X$ 上的非常数有理函数 $f$.因此，它到 $X_2$ 的限制全纯.因为 $X_2$ 紧，$f$ 在 $X_2$ 的每个连通分支上常值(Cartan, H., \emph{Elementary Theory of Analytic Functions of One or Several Variables}, Addison-Wesley, 1963, VI.4.5)，例如在某个无限连通分支上 $f(z) = a$.则 $f(z) - a$ 有无限多个零点，这与它是有理函数矛盾.

连通非奇异簇 $X$ 可以通过对维数作归纳证明它对复拓扑连通---见下面的注记~\ref{rem:21.9}(另见 Shafarevich, I., \emph{Basic Algebraic Geometry}, 1994, VII.2).

对 $r = 1$，定理断言 $X$ 与 $X^{\mathrm{an}}$ 的有限平展覆盖之间存在自然的一一对应.
\end{remark}

\begin{theorem}\label{thm:21.3}[Riemann 存在定理]
对 $\mathbb{C}$ 上的任何非奇异代数簇 $X$，函子 $Y \mapsto Y^{\mathrm{an}}$ 定义了 $X$ 与 $X^{\mathrm{an}}$ 的有限平展覆盖范畴之间的等价.
\end{theorem}

\begin{proof}
显然，这首先由 Riemann 对 Riemann 曲面证明.一般情形归功于 Grauert 与 Remmert.用奇点消解可以缩短证明.我概述 (SGA 1, XII) 的证明.

第 1 步：对任何射影非奇异代数簇 $X$，函子 $Y \mapsto Y^{\mathrm{an}}$ 定义了 $X$ 的有限覆盖范畴到 $X^{\mathrm{an}}$ 的有限覆盖范畴的等价.

事实上，对完备簇 $X$，函子 $M \mapsto M^{\mathrm{an}}$ 是 $X_{\mathrm{zar}}$ 上凝聚 $\mathcal{O}_X$-模范畴到 $X^{\mathrm{an}}$ 上凝聚 $\mathcal{O}_{X^{\mathrm{an}}}$-模范畴的等价(见 Serre, GAGA\footnote{\emph{G\'eom\'etrie alg\'ebrique et g\'eom\'etrie analytique}. Ann. Inst. Fourier, Grenoble 6 (1955--1956), 1--42.}, 1956).给出 $X$ 的有限覆盖就是给出凝聚 $\mathcal{O}_X$-模连同 $\mathcal{O}_X$-代数结构.由于 $X^{\mathrm{an}}$ 同理，论断显然.

第 2 步：从 $X$ 的有限覆盖范畴到 $X^{\mathrm{an}}$ 的有限覆盖范畴的函子 $Y \mapsto Y^{\mathrm{an}}$ 全忠实.

必须证明对 $X$ 的任何有限平展覆盖 $Y$ 与 $Y'$，
\[
\operatorname{Hom}_X(Y, Y') \xrightarrow{\sim} \operatorname{Hom}_{X^{\mathrm{an}}}(Y^{\mathrm{an}}, Y'^{\mathrm{an}}).
\]
可以假定 $X$ 连通.给出 $X$-态射 $Y \to Y'$ 就是给出 $Y \times_{X} Y' \to Y$ 的截面，而这等同于给出 $Y \times_{X} Y'$ 的连通分支 $\Gamma$，使得投影 $Y \times_{X} Y' \to Y$ 诱导的态射 $\Gamma \to X$ 是同构(见命题~\ref{prop:2.15}).但正如我们刚注意到的，$Y \times_{X} Y'$ 的连通分支与 $Y^{\mathrm{an}} \times_{X^{\mathrm{an}}} Y'^{\mathrm{an}}$ 的连通分支一致，且若 $\Gamma$ 是 $Y \times_{X} Y'$ 的连通分支，则投影 $\Gamma \to Y$ 是同构，当且仅当 $\Gamma^{\mathrm{an}} \to Y^{\mathrm{an}}$ 是同构.

还剩下证明函子 $Y \mapsto Y^{\mathrm{an}}$ 本质满.

第 3 步：问题对 $X$ 上的 Zariski 拓扑是局部的.

我们的意思是：若该函子对 $X$ 的某个 Zariski 开覆盖中的所有 $X_i$ 都本质满，则它对 $X$ 本身也本质满.这立即由第 2 步推出：由有限平展覆盖 $Y \to X^{\mathrm{an}}$，得到有限平展覆盖 $Y_i \to X^{\mathrm{an}}_i$ 与拼接数据 $Y_{ij} \to Y_{ji}$；若每个 $Y_i \to X^{\mathrm{an}}_i$ 都是代数的，则(由第 2 步)拼接数据也是代数的，从而给出 $X$ 的代数平展覆盖.

现在可以假定 $X$ 仿射.由奇点消解(Hironaka)，存在非奇异射影簇 $\widetilde{X}$ 与开浸入 $X \hookrightarrow \widetilde{X}$，把 $X$ 等同于 $\widetilde{X}$ 的稠密开子集，且 $\widetilde{X} \setminus X$ 是带法向相交的因子.在这些假设下，可以证明 $X^{\mathrm{an}}$ 的每个有限平展覆盖 $Y \to X^{\mathrm{an}}$ 都延拓为 $\widetilde{X}^{\mathrm{an}}$ 的有限覆盖(SGA 1, XII 5.3)，对它应用第 1 步.
\end{proof}

\begin{example}\label{ex:21.4}
最后一段需要 $Y$ 在 $X^{\mathrm{an}}$ 上平展的假设.设 $X$ 是复平面中的单位圆盘，$U$ 是 $X$ 中原点的补集，$U'$ 是由方程
\[
T^{2} = \sin\frac{1}{z}
\]
定义的 $U$ 的覆盖，其中 $z$ 是 $U$ 上的坐标函数.则 $U'$ 不延拓为 $X$ 的有限覆盖，因为若延拓，则 $X' \to X$ 不平展的点集是含所有满足 $\sin\frac{1}{z} = 0$ 的点的闭解析集，这是荒谬的.
\end{example}

从现在起假定 $X$ 连通.定理~\ref{thm:21.3} 蕴含对任何 $x \in X$，$\pi_1(X^{\mathrm{cx}}, x)$ 与 $\pi_1(X_{\mathrm{et}}, x)$ 有相同的有限商.因此，$X_{\mathrm{et}}$ 上带有限茎的局部常值层 $F$ 与 $X^{\mathrm{cx}}$ 上带有限茎的局部常值层之间存在自然的一一对应 $F \leftrightarrow F^{\mathrm{cx}}$.现在把定理重述为更强的形式：

\begin{theorem}\label{thm:21.5}
设 $X$ 是 $\mathbb{C}$ 上的连通非奇异簇.对 $X_{\mathrm{et}}$ 上带有限茎的任何局部常值层 $F$，对所有 $r \geq 0$，$H^{r}(X_{\mathrm{et}}, F) \cong H^{r}(X^{\mathrm{cx}}, F^{\mathrm{cx}})$.
\end{theorem}

从现在起，我省略 $F^{\mathrm{cx}}$ 上的上标.

考虑如下情形：$X$ 是带两种拓扑 $T_1$ 与 $T_2$(通常意义下)的集合，假定 $T_2$ 比 $T_1$ 细.设 $X_i$ 表示赋予拓扑 $T_i$ 的 $X$.因为 $T_2$ 比 $T_1$ 细，恒等映射 $f \colon X_2 \to X_1$ 连续.我们问：对 $X_2$ 上的层 $F$，什么时候对所有 $r \geq 0$，$H^{r}(X_1, f_{*}F) \cong H^{r}(X_2, F)$？注意 $f_{*}F$ 无非是 $F$ 到 $X_1$ 的限制.Leray 谱序列
\[
H^{r}(X_1, R^{s}f_{*}F) \Rightarrow H^{r+s}(X_2, F)
\]
给出了答案.即当 $s > 0$ 时若 $R^{s}f_{*}F = 0$，上同调群就一致.但 $R^{s}f_{*}F$ 是与预层 $U \mapsto H^{s}(U_2, F)$ 相伴的层(这里 $U$ 是 $X_1$ 的开子集，$U_2$ 表示赋予 $T_2$-拓扑的同一集合).于是得到 $H^{r}(X_1, f_{*}F) \cong H^{r}(X_2, F)$ 的如下判别准则：
\begin{quote}
对任何开集 $U \subset X_1$ 与任何 $t \in H^{s}(U_2, F)$($s > 0$)，存在覆盖 $U = \bigcup U^{(i)}$(对 $T_1$-拓扑)，使得 $t$ 在每个 $H^{s}(U^{(i)}_2, F)$ 中映到零.
\end{quote}
松散地说，如果 $T_1$-拓扑细到足以消灭 $T_2$-上同调类，我们就说它细到足以计算 $F$ 的 $T_2$-上同调.

类似的评注适用于景的连续映射.设 $X^{\mathrm{ecx}}$ 表示赋予 Grothendieck 拓扑的 $X$，其中覆盖是 $X$ 上复流形的平展映射的满族 $(U_i \to U)$.有(景的)连续态射
\[
X^{\mathrm{cx}} \to X^{\mathrm{ecx}} \to X_{\mathrm{et}}.
\]
逆映射定理表明，复流形 $U$ 的每个复-平展覆盖 $(U_i \to U)$ 都有开覆盖(通常意义下)的加细.因此，左边的箭头给出上同调的同构.还剩下证明右边的箭头也如此，为此，上面的讨论表明只需证明如下论断：

\begin{lemma}\label{lem:21.6}
设 $U$ 是连通非奇异簇，$F$ 是 $U^{\mathrm{ecx}}$ 上带有限茎的局部常值层.对任何 $t \in H^{s}(U^{\mathrm{cx}}, F)$($s > 0$)，存在平展覆盖 $U_i \to U$(代数意义下)，使得 $t$ 在每个 $H^{s}(U_{i,\mathrm{cx}}, F)$ 中映到零.
\end{lemma}

引理证明的思路如下.对 $U$ 的维数作归纳.显然问题对 $U$ 上的平展拓扑是局部的，即只需对 $U$ 的某个平展覆盖 $(U_i \to U)_{i \in I}$ 中的每个 $U_i$，证明 $t$ 的像 $t_i \in H^{r}(U_i, F)$ 的论断.于是可以假定 $F$ 常值，且 $U$ 已换成某个``Zariski-小''集 $U$.然后找到态射 $f \colon U \to S$(从 $U$ 到较低维数的非奇异簇 $S$)，使得对 $s \neq 0, 1$，$R^{s}f_{*}F$ 为零，对 $s = 0, 1$ 局部常值(对复拓扑的直像).$f$ 的 Leray 谱序列给出正合序列
\[
\cdots \to H^{s}(S^{\mathrm{cx}}, f_{*}F) \to H^{s}(U^{\mathrm{cx}}, F) \to H^{s-1}(S^{\mathrm{cx}}, R^{1}f_{*}F) \to \cdots.
\]
设 $t$ 映到 $H^{s-1}(S^{\mathrm{cx}}, R^{1}f_{*}F)$ 中的 $t''$.若 $s - 1 > 0$，则由归纳，存在平展覆盖 $S_i \to S$(代数意义下)，使得 $t''$ 在每个 $H^{s-1}(S^{\mathrm{cx}}_i, R^{1}f_{*}F)$ 中限制为零.把 $U$ 换成 $U \times_{S} S_i$ 后，可以假定 $t$ 是 $H^{r}(S^{\mathrm{cx}}, f_{*}F)$ 中某个元素 $t'$ 的像，并再次应用归纳.这完成了 $s > 1$ 情形的证明.$s = 1$ 的情形由 Riemann 存在定理推出.

\begin{definition}\label{def:21.7}
初等纤维化是簇的正则映射 $f \colon U \to S$，它可以嵌入交换图
\[
\begin{tikzcd}[row sep=small, column sep=small]
U \arrow[r, "j"] \arrow[rd, "f"'] & Y \arrow[d, "h"] & Z \arrow[l, "i"'] \arrow[ld, "g"] \\
& S &
\end{tikzcd}
\]
其中：
\begin{enumerate}
  \item $j$ 是开浸入，$j(U)$ 在 $h$ 的每个纤维中稠密，且 $Y = i(Z) \cup j(U)$；
  \item $h$ 光滑且射影，纤维几何不可约且维数为 1；
  \item $g$ 有限且平展，$g$ 的每个纤维非空.
\end{enumerate}
\end{definition}

\begin{proposition}\label{prop:21.8}
设 $X$ 是代数闭域 $k$ 上的非奇异簇.对 $X$ 的任何闭点 $x$，存在初等纤维化 $U \to S$，其中 $U$ 是 $x$ 的开邻域，$S$ 非奇异.
\end{proposition}

\begin{proof}
思路是找到嵌入 $U \hookrightarrow \mathbb{P}^{r}$，使得 $U$ 的闭包 $\bar{U}$ 正规，并且存在特别好的投影映射 $\mathbb{P}^{r} \setminus E \to \mathbb{P}^{m-1}$($m = \dim U$).在投影映射的中心爆破 $U$ 后，得到初等纤维化 $U \to S$.见 SGA 4, XI(也许某天我会把它放进 AG).
\end{proof}

\begin{remark}\label{rem:21.9}
用该命题可以容易地完成定理~\ref{thm:21.1} 对 $r = 0$ 的证明.由于去掉或加上实余维数 $\geq 2$ 的分支不会连接或断开流形，可以把 $X$ 换成命题中的开子集 $U$，再换成 $Y$.于是，必须证明 $Y$ 对复拓扑连通(假定 $S$ 连通).设 $Y = Y_1 \sqcup Y_2$，$Y_1$ 与 $Y_2$ 都开且闭.因为 $h$ 的纤维连通，$Y_1$ 与 $Y_2$ 的每个都必须是纤维的并，即 $Y_i = h^{-1}h(Y_i)$.因为 $h$ 紧合，$h(Y_1)$ 与 $h(Y_2)$ 闭，因此除非一个是空的，它们将断开 $S$.

现在我们证明引理~\ref{lem:21.6}.因为论断对 $U$ 上的平展拓扑是局部的，可以假定 $U$ 有初等纤维化且 $F$ 常值，$F = \Lambda$.还有解析拓扑的 Gysin 正合序列
\[
\cdots \to H^{r-2}(Z^{\mathrm{cx}}, \Lambda(-1)) \to H^{r}(Y^{\mathrm{cx}}, \Lambda) \to H^{r}(U^{\mathrm{cx}}, \Lambda) \to \cdots.
\]
对任何复-开子集 $S' \subset S$，得到类似序列，只把 $Y, Z, U$ 换成 $h^{-1}(S')$，$g^{-1}(S')$,$f^{-1}(U)$.这样，得到 $S^{\mathrm{cx}}$ 上预层的正合序列，相伴的层序列是
\[
\cdots \to R^{r-2}g_{*}\Lambda(-1) \to R^{r}h_{*}\Lambda \to R^{r}f_{*}\Lambda \to \cdots
\]
(对复拓扑的高阶直像).由此推出层 $R^{r}f_{*}\Lambda$ 局部常值，且对所有 $s \in S$，自然映射 $(R^{r}f_{*}\Lambda)_s \to H^{r}(U_s, \Lambda)$ 是同构(把拓扑的紧合基变换定理应用于 $g$ 与 $h$，并用五引理).因此对 $r > 1$，$R^{r}f_{*}\Lambda = 0$.(\ref{def:21.7}) 之前的讨论表明这完成了定理~\ref{thm:21.5} 的证明.
\end{remark}

\begin{remark}\label{rem:21.10}
我们现在知道，对特征 $p \neq 0$ 的代数闭域 $k$ 上可提升为特征零的完备非奇异簇 $X_1$ 的任何完备非奇异簇 $X_0$，对所有 $r$，$H^{r}(X_0, \Lambda) \cong H^{r}(X^{\mathrm{cx}}_1, \Lambda)$.
\end{remark}
